% =====================================================================
% La frontera de medición: M1 y M3
% Parte I. El dato
%
% ENCARGO: dos coberturas, dos preguntas. La bifurcacion es en si misma
% un resultado: invierte los papeles de los agentes.
%
% Reglas del documento:
% - Una idea = una subseccion = una figura.
% - Toda figura declara su procedencia con \fuente{...}.
% - Toda cifra sale del canon de agosto; ver GUIA.md.
% - M1 y M3 se reportan siempre en paralelo.
% - Ningun superlativo sin releer el CSV de la frontera correcta.
% =====================================================================
\section{La frontera de medición: M1 y M3}
\label{sec:frontera}
\justifying

El capítulo anterior terminó con dos matrices y una advertencia: la
bifurcación entre las coberturas M1 y M3 ocurre en la primera etapa del
pipeline y en ninguna otra. Este capítulo desarrolla esa bifurcación, y lo
hace en capítulo aparte por una razón que conviene anticipar: no se trata
de dos versiones del mismo cálculo, ni de un análisis de sensibilidad
sobre un parámetro. Se trata de dos preguntas distintas, y las dos tienen
respuesta. Elegir una sola y presentarla como \emph{el} resultado sería la
decisión más discutible que este trabajo podría tomar.

\subsection{Dos preguntas, no dos versiones}
\label{sub:frontera-preguntas}

Cada institución de la comunidad tiene más de un medidor, y el
inventario de instalación del Capítulo~\ref{sec:dato} estableció qué
circuito registra cada uno: el Medidor 1, el circuito principal o de
inyección; el Medidor 3, un circuito secundario. Ninguno totaliza la
institución completa ni aísla el ramal que alimenta el sistema
fotovoltaico. Leer uno u otro no es una preferencia técnica:
define contra qué se compara la generación, y por tanto qué pregunta se
está respondiendo. La Figura~\ref{fig:frontera-esquema} ubica las dos
fronteras sobre el circuito de una institución.

\begin{figure}[H]
 \centering
 \begin{tikzpicture}[
 font=\scriptsize,
 blk/.style={draw=black!55, fill=black!3, rounded corners=1.5pt,
 minimum width=1.75cm, minimum height=0.6cm,
 align=center},
 med/.style={circle, draw=black!70, fill=white, minimum size=0.62cm,
 inner sep=0pt, font=\scriptsize\bfseries},
 inv/.style={draw=successcolor, fill=successcolor!14, very thick,
 rounded corners=1.5pt, minimum width=1.35cm,
 minimum height=0.55cm, align=center},
 lin/.style={-{Stealth[length=3.5pt]}, draw=black!60},
 et/.style={font=\scriptsize\itshape, align=left},
 ]
 \node[blk] (red) at (0,0) {Red};
 \node[med, draw=sectioncolor, fill=sectioncolor!14, very thick]
 (m1) at (2.0,0) {M1};
 \node[blk, minimum width=2.1cm] (campus) at (4.4,0)
 {Circuito principal\\o de inyección};
 \node[inv] (inv) at (7.1,0) {Inversor};

 \node[med, draw=warningcolor, fill=warningcolor!16, very thick]
 (m3) at (4.4,-1.35) {M3};
 \node[blk, minimum width=2.1cm] (circ) at (4.4,-2.55)
 {Circuito\\secundario};

 \draw[lin] (red) -- (m1);
 \draw[lin] (m1) -- (campus);
 \draw[lin] (inv) -- (campus);
 \draw[lin] (campus) -- (m3);
 \draw[lin] (m3) -- (circ);

 \node[et, text=sectioncolor, right=0.5cm of inv, anchor=west,
 text width=3.9cm]
 {\textbf{M1} mide el circuito principal\\
 o de inyección: la generación\\
 equivale al 19,1 \% de su consumo};
 \node[et, text=warningcolor, right=0.5cm of circ, anchor=west,
 text width=3.9cm]
 {\textbf{M3} mide un circuito\\
 secundario: la generación\\
 equivale al 91,2 \% de su consumo};
 \end{tikzpicture}
 \caption[Las dos fronteras de medición]{Posición de las dos fronteras
 de medición en el circuito de una institución. El mismo inversor
 entrega la misma energía en los dos casos; lo que cambia es la carga
 contra la que se la compara. El esquema es una idealización: la posición
 del punto de inyección respecto de cada medidor varía entre
 instituciones, como recoge el inventario del Capítulo~\ref{sec:dato}.
 \fuente{elaboración propia a partir de \code{DEMAND\_METER\_CONFIG} y
 \code{PAPER\_METER\_DEMAND\_CONFIG} en \file{data/preprocessing.py}.}}
 \label{fig:frontera-esquema}
\end{figure}

La cobertura \textbf{M1} pregunta: \emph{¿qué le aporta este mercado a
una carga grande, frente a la cual la generación es escasa?} La cobertura
\textbf{M3} pregunta: \emph{¿qué le aporta a una carga del orden de la
propia generación?} La primera describe un mercado con poco excedente que
repartir; la segunda, uno donde el excedente abunda. Ninguna es más
legítima que la otra, porque la relación de cada circuito con el consumo
institucional completo está por establecer.

\subsection{Qué mide cada frontera}
\label{sub:frontera-cobertura}

La diferencia entre las dos coberturas se resume en un hecho que conviene
subrayar porque desactiva de entrada una posible objeción: \textbf{la
generación es idéntica en las dos}. El pipeline designa el mismo inversor
por institución en un caso y en el otro, de modo que la matriz $G$ es la
misma matriz. Lo único que cambia es el denominador. La
Figura~\ref{fig:frontera-cobertura} confronta las dos energías en el
agregado comunitario y en cada institución.

\begin{figure}[H]
 \centering
 \includegraphics[width=0.95\textwidth]{figuras/f4_02_cobertura.png}
 \caption[Cobertura de generación en cada frontera]{Energía medida en
 el horizonte bajo cada frontera. A la izquierda, la comunidad
 completa: la barra de generación es la misma en los dos casos y solo
 cambia la de demanda. A la derecha, la razón entre generación y
 demanda por institución. Las cifras se leen del caché del
 preprocesamiento, no se fijan en el código.
 \fuente{elaboración propia desde el caché de estados intermedios del
 preprocesamiento.}}
 \label{fig:frontera-cobertura}
\end{figure}

\notafig{La demanda comunitaria medida es de \uni{318046}{kWh} bajo M1
y de \uni{66734}{kWh} bajo M3, mientras que la generación es de
\uni{60860}{kWh} en ambos casos, es decir, exactamente la misma cifra. De
ahí las dos coberturas, \pct{19.1} y \pct{91.2}, que no reflejan dos
instalaciones distintas sino dos maneras de delimitar la carga.
Elaboración propia desde el caché de estados intermedios.}

\begin{lecturabox}
Conviene resistir la lectura de que M3 «tiene más generación». No la
tiene: tiene la misma. Lo que M3 hace es comparar esa generación contra
una carga cinco veces menor, porque solo mira un circuito secundario de
cada institución. Una generación que equivale al \pct{19} del consumo del
circuito principal puede equivaler al \pct{91} del consumo de un circuito
menor sin que se haya instalado un solo panel adicional.
\end{lecturabox}

\subsection{Los mismos días, dos lecturas}
\label{sub:frontera-perfiles}

El mismo hecho visto en el eje del tiempo resulta más elocuente que en
totales. La Figura~\ref{fig:frontera-perfiles} reúne el perfil diario de
las dos coberturas en un mismo par de paneles.

\begin{figure}[H]
 \centering
 \includegraphics[width=0.95\textwidth]{figuras/f4_03_perfiles_m1_m3.png}
 \caption[Perfil comunitario bajo cada frontera]{Perfil medio por hora
 del día de la comunidad en cada cobertura. La curva verde es la misma
 en los dos paneles; la gris cae casi un factor cinco. En M3 las dos
 curvas llegan a cruzarse en las horas centrales, lo que en M1 no
 ocurre nunca.
 \fuente{elaboración propia desde el caché de estados intermedios del
 preprocesamiento.}}
 \label{fig:frontera-perfiles}
\end{figure}

\notafig{El cruce del panel derecho es la consecuencia operativa de
todo el capítulo. Bajo M1 la generación no alcanza a la demanda
comunitaria en ninguna hora del día, de modo que la comunidad es siempre
compradora neta frente a la red y el intercambio interno se limita a
redistribuir un excedente escaso. Bajo M3 hay horas en que la comunidad
genera más de lo que su circuito fotovoltaico consume, y el excedente
disponible para intercambio es sustancialmente mayor. Elaboración propia
desde el caché de estados intermedios.}

\subsection{La inversión de los papeles}
\label{sub:frontera-inversion}

Hasta aquí el efecto es de escala, y podría parecer que M3 es simplemente
«M1 con más recurso relativo». No lo es. El cambio de denominador altera
qué institución tiene excedente y cuál tiene déficit en cada hora, y el
resultado es que los papeles se invierten. La
Figura~\ref{fig:frontera-inversion} cruza las dos coberturas conservando
el orden de las instituciones, de modo que el vuelco se lee de un panel a
otro.

\begin{figure}[H]
 \centering
 \includegraphics[width=0.95\textwidth]{figuras/f4_04_inversion_papeles.png}
 \caption[Reparto de la energía transada en cada frontera]{Participación
 de cada institución en la energía transada, hacia la derecha como
 vendedora y hacia la izquierda como compradora. Las cinco conservan el
 mismo orden vertical en los dos paneles, de modo que el cambio de
 frontera se lee como un vuelco del perfil.
 \fuente{elaboración propia desde
 \file{p2p\_breakdown\_flujos.csv} de las dos carpetas canónicas.}}
 \label{fig:frontera-inversion}
\end{figure}

\notafig{Bajo M1, Udenar aporta el \pct{76.3} de la energía vendida y
la UCC absorbe el \pct{36.3} de la comprada. Bajo M3, Udenar baja al
\pct{30.5} de lo vendido, el HUDN pasa al primer lugar con el \pct{33.9},
y Cesmag, que en M1 vendía el \pct{8.0}, absorbe el \pct{70.9} de todo lo
comprado. No es un reordenamiento menor: es un cambio de papel.
Elaboración propia desde \file{p2p\_breakdown\_flujos.csv} de las dos
carpetas canónicas.}

\begin{trampabox}
Esta figura es el origen de la regla de redacción más estricta del
documento. Ninguna frase con superlativo («el mayor vendedor», «solo dos
instituciones», «ninguna») puede escribirse sin decir de qué cobertura
habla, porque la institución que domina la venta en una frontera apenas
participa en la otra. Toda afirmación de ese tipo en los capítulos
siguientes se comprobó releyendo el archivo de flujos de la frontera
correspondiente.
\end{trampabox}

\subsection{El embudo de horas de mercado}
\label{sub:frontera-embudo}

Queda una magnitud que ninguna de las figuras anteriores muestra y que
conviene fijar antes de llegar a los resultados: de las \num{6144} horas
del horizonte, en muy pocas llega a haber mercado. La
Figura~\ref{fig:frontera-embudo} presenta esa reducción en cuatro
escalones sucesivos.

\begin{figure}[H]
 \centering
 \includegraphics[width=0.95\textwidth]{figuras/f4_05_embudo_horas.png}
 \caption[Del horizonte a las horas con mercado]{Reducción sucesiva del
 número de horas: horizonte completo, horas con generación, horas en
 que alguna institución tiene excedente sobre su propia demanda, y
 horas en que efectivamente se liquidó un intercambio. El último
 escalón procede de los archivos de flujos y no de una estimación.
 \fuente{elaboración propia desde el caché del preprocesamiento y
 \file{p2p\_breakdown\_flujos.csv} de las dos carpetas canónicas.}}
 \label{fig:frontera-embudo}
\end{figure}

\notafig{El recorrido es el mismo en las dos coberturas hasta el
segundo escalón, porque hay generación en \num{3353} horas y esa cifra no
depende de la frontera. La separación empieza en el tercero: bajo M1
quedan \num{1047} horas con excedente y se transa en \num{1029}; bajo M3,
\num{2512} y \num{1752}. Elaboración propia desde el caché del
preprocesamiento y los archivos de flujos.}

\begin{lecturabox}
Que el mercado opere en el \pct{16.7} de las horas bajo M1 y en el
\pct{28.5} bajo M3 no es un defecto del mecanismo: es una consecuencia
del recurso. Un mercado de excedentes solares no puede funcionar de noche.
Lo que la figura acota es la ambición razonable del ejercicio: cualquier
ventaja que el Capítulo~\ref{sec:res} reporte se genera en esa fracción de
horas y se diluye sobre el total.
\end{lecturabox}

\subsection{Una salvedad que hay que declarar}
\label{sub:frontera-mariana}

La cobertura M3 no es perfectamente homogénea, y conviene decirlo aquí y
no en una nota al pie. Cuatro de las cinco instituciones tienen un
medidor de circuito secundario que el pipeline lee directamente. Mariana
no lo tiene. Para representarla en M3, el pipeline
reutiliza su medidor de circuito principal escalado por un factor de $0{,}3$, bajo
el supuesto de que el circuito relevante equivale aproximadamente a una
facultad.

Ese factor es una convención, no una medición, y se declara como tal.
Afecta a una de las cinco instituciones y solo en una de las dos
coberturas. Su efecto sobre el resultado agregado se examina en el
Capítulo~\ref{sec:rob}, donde el barrido del nivel de demanda cubre
holgadamente el rango en que ese factor podría equivocarse.

\subsection{La consecuencia: reportar dos veces}
\label{sub:frontera-consecuencia}

De todo lo anterior se sigue la decisión metodológica que gobierna el
resto del documento: \textbf{cada resultado se reporta bajo las dos
coberturas, siempre, y sin elegir una como principal.}

La tentación de elegir es real y conviene nombrarla. M1 es la frontera en
la que el mercado entre pares obtiene su mejor resultado relativo, y una
presentación que se quedara solo con M1 mostraría un trabajo más
concluyente. M3 es la frontera en la que la generación pesa más sobre
la carga medida, con una razón entre generación y demanda del
\pct{91.2}; y es también la frontera en la que, como se verá, el
mecanismo colectivo supera el mercado entre pares. Presentar solo una
de las dos sería, en el primer caso, seleccionar el resultado favorable, y
en el segundo, ocultar el caso en que el mecanismo propuesto funciona
mejor.

La convención adoptada evita las dos: las figuras de los capítulos de
resultados llevan siempre dos paneles, el izquierdo M1 y el derecho M3, y
cuando las dos coberturas dan lecturas opuestas, se dicen las dos y se
discute por qué difieren.

En síntesis, la bifurcación deja dos preguntas con respuesta y una regla
de reporte: dos paneles por figura, con una razón entre generación y
demanda del \pct{19.1} bajo M1 y del \pct{91.2} bajo M3. Con la energía
así delimitada queda pendiente la otra mitad del dato, los precios, que el
capítulo siguiente construye desde sus fuentes primarias.
